@thesis{al-aukaily2017response,
  type = {phdthesis},
  title = {Response {{Sensitivity Formulations}} for {{Geometrically Nonlinear Finite}}  {{Element Analysis}}},
  author = {Al-Aukaily, Alabbas},
  date = {2017},
  institution = {Oregon State University}
}
% == BibLateX quality report for al-aukaily2017response:
% ? Title looks like it was stored in title-case in Zotero
% ? unused Type ("Ph. D.")

@article{aln2015fenics,
  title = {The {{FEniCS Project Version}} 1.5},
  author = {Aln, Martin S and Kehlet, Benjamin and Logg, Anders and Richardson, Chris and Ring, Johannes and Rognes, E and Wells, Garth N},
  date = {2015},
  doi = {10.11588/ans.2015.100.20553},
  abstract = {The FEniCS Project is a collaborative project for the development of innovative concepts and tools for automated scientific computing, with a particular focus on the solution of differential equations by finite element methods. The FEniCS Project consists of a collection of interoperable software components, including DOLFIN, FFC, FIAT, Instant, mshr, and UFL. This note describes the new features and changes introduced in the release of FEniCS version 1.5.},
  langid = {english}
}
% == BibLateX quality report for aln2015fenics:
% Missing required field 'journaltitle'
% ? Title looks like it was stored in title-case in Zotero
% ? unused Library catalog ("Zotero")

@software{autodiff,
  title = {Autodiff},
  date = {2021-04-16T10:11:06Z},
  origdate = {2018-07-19T07:28:38Z},
  url = {https://github.com/autodiff/autodiff},
  urldate = {2021-04-17},
  abstract = {automatic differentiation made easier for C++. Contribute to autodiff/autodiff development by creating an account on GitHub.},
  organization = {autodiff}
}
% == BibLateX quality report for autodiff:
% ? Title looks like it was stored in lower-case in Zotero
% ? unused Library catalog ("GitHub")
% ? unused Programming language ("C++")

@article{bangerth2007deal,
  title = {Deal.{{II---A}} General-Purpose Object-Oriented Finite Element Library},
  author = {Bangerth, W. and Hartmann, R. and Kanschat, G.},
  date = {2007-08-01},
  journaltitle = {ACM Transactions on Mathematical Software},
  shortjournal = {ACM Trans. Math. Softw.},
  volume = {33},
  number = {4},
  pages = {24-es},
  issn = {00983500},
  doi = {10.1145/1268776.1268779},
  langid = {english},
  annotation = {813 citations (Crossref) [2022-06-27]}
}
% == BibLateX quality report for bangerth2007deal:
% ? unused Library catalog ("DOI.org (Crossref)")
% ? unused Url ("http://portal.acm.org/citation.cfm?doid=1268776.1268779")

@article{barthelemy1995automatic,
  title = {Automatic Differentiation as a Tool in Engineering Design},
  author = {Barthelemy, J. F. and Hall, L. E.},
  date = {1995},
  doi = {10.1007/BF01758823},
  abstract = {Automatic Differentiation (AD) is a tool that systematically implements the chain rule of differentiation to obtain the derivatives of functions calculated by computer programs. In this paper, it is assessed as a tool for engineering design. The paper discusses the forward and reverse modes of AD, their computing requirements, as well as approaches to implementing AD. It continues with the application of two different tools to two medium-size structural analysis problems to generate sensitivity information typically necessary in an optimization or design situation. The paper concludes with the observation that AD is to be preferred to finite differencing in most cases, as long as sufficient computer storage is available; in some instances, AD may be the alternative to consider in lieu of analytical sensitivity analysis.},
  langid = {english},
  annotation = {17 citations (Crossref) [2022-06-27]}
}
% == BibLateX quality report for barthelemy1995automatic:
% Missing required field 'journaltitle'
% ? unused Library catalog ("Zotero")

@article{bartholomew-biggs2000automatic,
  title = {Automatic Differentiation of Algorithms},
  author = {Bartholomew-Biggs, Michael and Brown, Steven and Christianson, Bruce and Dixon, Laurence},
  date = {2000-12-01},
  journaltitle = {Journal of Computational and Applied Mathematics},
  shortjournal = {Journal of Computational and Applied Mathematics},
  series = {Numerical {{Analysis}} 2000. {{Vol}}. {{IV}}: {{Optimization}} and {{Nonlinear Equations}}},
  volume = {124},
  number = {1},
  pages = {171--190},
  issn = {0377-0427},
  doi = {10/b94pb4},
  abstract = {We introduce the basic notions of automatic differentiation, describe some extensions which are of interest in the context of nonlinear optimization and give some illustrative examples.},
  langid = {english},
  annotation = {106 citations (Crossref) [2022-06-27]}
}
% == BibLateX quality report for bartholomew-biggs2000automatic:
% ? unused Library catalog ("ScienceDirect")
% ? unused Url ("https://www.sciencedirect.com/science/article/pii/S0377042700004222")

@article{baydin2018automatic,
  title = {Automatic {{Differentiation}} in {{Machine Learning}}: A {{Survey}}},
  author = {Baydin, Atılım Gunes and Pearlmutter, Barak A and Radul, Alexey Andreyevich and Siskind, Jeffrey Mark},
  date = {2018},
  pages = {43},
  abstract = {Derivatives, mostly in the form of gradients and Hessians, are ubiquitous in machine learning. Automatic differentiation (AD), also called algorithmic differentiation or simply “autodiff”, is a family of techniques similar to but more general than backpropagation for efficiently and accurately evaluating derivatives of numeric functions expressed as computer programs. AD is a small but established field with applications in areas including computational fluid dynamics, atmospheric sciences, and engineering design optimization. Until very recently, the fields of machine learning and AD have largely been unaware of each other and, in some cases, have independently discovered each other’s results. Despite its relevance, general-purpose AD has been missing from the machine learning toolbox, a situation slowly changing with its ongoing adoption under the names “dynamic computational graphs” and “differentiable programming”. We survey the intersection of AD and machine learning, cover applications where AD has direct relevance, and address the main implementation techniques. By precisely defining the main differentiation techniques and their interrelationships, we aim to bring clarity to the usage of the terms “autodiff”, “automatic differentiation”, and “symbolic differentiation” as these are encountered more and more in machine learning settings.},
  langid = {english}
}
% == BibLateX quality report for baydin2018automatic:
% Missing required field 'journaltitle'
% ? unused Library catalog ("Zotero")

@article{beck1994automatic,
  title = {Automatic Differentiation of Iterative Processes},
  author = {Beck, Thomas},
  date = {1994-05},
  journaltitle = {Journal of Computational and Applied Mathematics},
  shortjournal = {Journal of Computational and Applied Mathematics},
  volume = {50},
  number = {1-3},
  pages = {109--118},
  issn = {03770427},
  doi = {10.1016/0377-0427(94)90293-3},
  abstract = {Automatic differentiation is used to compute the values of the derivative of a function. If the function is given by a computational graph or code list, then the derivative values can be obtained using the chain rule. An iterative process can be regarded as an infinite code list. It is well known from classical analysis that the limit of the derivatives of the code list is not necessarily equal to the derivative of the limit function. The limit of the derivatives is correct for an important class of iterative processes including generalized Newton methods.},
  langid = {english},
  annotation = {9 citations (Crossref) [2022-06-27]}
}
% == BibLateX quality report for beck1994automatic:
% ? unused Library catalog ("DOI.org (Crossref)")
% ? unused Url ("https://linkinghub.elsevier.com/retrieve/pii/0377042794902933")

@software{clad,
  title = {Clad},
  author = {Vassilev, Vassil},
  date = {2021-04-14T22:47:30Z},
  origdate = {2014-02-21T08:45:02Z},
  url = {https://github.com/vgvassilev/clad},
  urldate = {2021-04-18},
  abstract = {clad -- automatic differentiation for C/C++. Contribute to vgvassilev/clad development by creating an account on GitHub.}
}
% == BibLateX quality report for clad:
% ? Title looks like it was stored in lower-case in Zotero
% ? unused Library catalog ("GitHub")
% ? unused Programming language ("C++")

@inproceedings{clough1960finite,
  title = {The Finite Element Method in Plane Stress Analysis},
  author = {Clough, Ray W.},
  date = {1960}
}
% == BibLateX quality report for clough1960finite:
% Missing required field 'booktitle'

@article{dadvand2010objectoriented,
  ids = {dadvandObjectorientedEnvironmentDeveloping2010a},
  title = {An {{Object-oriented Environment}} for {{Developing Finite Element Codes}} for {{Multi-disciplinary Applications}}},
  author = {Dadvand, Pooyan and Rossi, Riccardo and Oñate, Eugenio},
  date = {2010-09},
  journaltitle = {Archives of Computational Methods in Engineering},
  shortjournal = {Arch Computat Methods Eng},
  volume = {17},
  number = {3},
  pages = {253--297},
  issn = {1134-3060, 1886-1784},
  doi = {10.1007/s11831-010-9045-2},
  abstract = {The objective of this work is to describe the design and implementation of a framework for building multidisciplinary finite element programs. The main goals are generality, reusability, extendibility, good performance and memory efficiency. Another objective is preparing the code structure for team development to ensure the easy collaboration of experts in different fields in the development of multi-disciplinary applications.},
  langid = {english},
  annotation = {203 citations (Crossref) [2022-06-27]}
}
% == BibLateX quality report for dadvand2010objectoriented:
% 'issn': not a valid ISSN
% ? unused Library catalog ("DOI.org (Crossref)")
% ? unused Url ("http://link.springer.com/10.1007/s11831-010-9045-2")

@software{DEALii,
  title = {Dealii/Dealii},
  date = {1997},
  origdate = {2014-07-17T14:00:26Z},
  location = {College Station, Tx},
  url = {https://github.com/dealii/dealii},
  urldate = {2019-12-02},
  abstract = {Comprehensive set of tools for finite element codes, scaling from laptops to clusters with 10,000+ cores. Written in C++.},
  organization = {Texas A\&M University},
  version = {1.0}
}
% == BibLateX quality report for DEALii:
% ? Title looks like it was stored in lower-case in Zotero
% ? unused Library catalog ("GitHub")
% ? unused Programming language ("C++")

@article{derkiureghian2006structural,
  title = {Structural Reliability Software at the {{University}} of {{California}}, {{Berkeley}}},
  author = {Der Kiureghian, Armen and Haukaas, T. and Fujimura, K.},
  date = {2006-01},
  journaltitle = {Structural Safety},
  shortjournal = {Structural Safety},
  volume = {28},
  number = {1-2},
  pages = {44--67},
  issn = {01674730},
  doi = {10.1016/j.strusafe.2005.03.002},
  abstract = {The paper provides an overview of three software for structural reliability analysis developed at the University of California, Berkeley. These are CalREL, which is a general-purpose structural reliability code written in FORTRAN; FERUM, which is a collection of MatlabÒ files that can be used for structural reliability analysis in conjunction with simple finite element models; and OpenSees, which is an object-oriented code for nonlinear structural response simulation with reliability capabilities. Reliability analysis methods featured by these software include FORM, SORM, Monte Carlo simulation, and various importance sampling methods. In addition, variable importance and sensitivity measures are provided.},
  langid = {english},
  annotation = {78 citations (Crossref) [2022-06-27]}
}
% == BibLateX quality report for derkiureghian2006structural:
% ? unused Library catalog ("DOI.org (Crossref)")
% ? unused Url ("https://linkinghub.elsevier.com/retrieve/pii/S0167473005000160")

@software{elle-0050,
  title = {Parameterized Finite Element Analysis of a {{Poisson}} Problem in {{2D}}},
  author = {Perez, Claudio},
  date = {2021-05},
  url = {https://claudioperez.github.io/anabel/gallery/elle-0050/},
  organization = {University of California, Berkeley}
}
% == BibLateX quality report for elle-0050:
% ? unused System ("Anabel")

@article{elliott2009beautiful,
  ids = {elliottBeautifulDifferentiation2009},
  title = {Beautiful {{Differentiation}}},
  author = {Elliott, Conal M},
  date = {2009-08-31},
  journaltitle = {ACM SIGPLAN Notices},
  shortjournal = {SIGPLAN Not.},
  volume = {44},
  number = {9},
  pages = {12},
  issn = {03621340},
  doi = {10.1145/1631687.1596579},
  abstract = {Automatic differentiation (AD) is a precise, efficient, and convenient method for computing derivatives of functions. Its forwardmode implementation can be quite simple even when extended to compute all of the higher-order derivatives as well. The higherdimensional case has also been tackled, though with extra complexity. This paper develops an implementation of higher-dimensional, higher-order, forward-mode AD in the extremely general and elegant setting of calculus on manifolds and derives that implementation from a simple and precise specification.},
  langid = {english},
  annotation = {1 citations (Crossref) [2022-06-27]}
}
% == BibLateX quality report for elliott2009beautiful:
% ? Title looks like it was stored in title-case in Zotero
% ? unused Library catalog ("DOI.org (Crossref)")
% ? unused Url ("http://portal.acm.org/citation.cfm?doid=1631687.1596579")

@software{feap,
  title = {{{FEAP}}},
  url = {http://feap.berkeley.edu/},
  key = {FEAP}
}
% == BibLateX quality report for feap:
% Unexpected field 'key'
% ? unused Programming language ("Fortran")

@software{FEDEASLab,
  title = {{{FEDEASLab}}},
  author = {Filippou, Filip C.},
  date = {1998},
  location = {Berkeley, Ca},
  organization = {University of California, Berkeley}
}
% == BibLateX quality report for FEDEASLab:
% ? unused Programming language ("Matlab")

@software{FEDEASLab,
  title = {{{FEDEASLab}}}
}

@software{FEniCS,
  title = {{{FEniCS Project}}},
  url = {https://fenicsproject.org/},
  urldate = {2019-11-22},
  abstract = {FEniCS is a popular open-source (LGPLv3) computing platform for solving partial differential equations (PDEs). FEniCS enables users to quickly translate scientific models into efficient finite element code. With the high-level Python and C++ interfaces to FEniCS, it is easy to get started, but FEniCS offers also powerful capabilities for more experienced programmers. FEniCS runs on a multitude of platforms ranging from laptops to high-performance clusters.},
  annotation = {github: https://github.com/FEniCS}
}
% == BibLateX quality report for FEniCS:
% ? Title looks like it was stored in title-case in Zotero
% ? unused Programming language ("Python, C++")

@report{filippou2004fedeaslab,
  title = {{{FEDEASLab}}: {{A Matlab Toolbox}} for {{Nonlinear Structural Response Simulations}}},
  author = {Filippou, Filip C.},
  date = {2004},
  number = {UCB/SEMM 2004/01},
  institution = {University of California Berkeley},
  location = {Berkeley, Ca}
}
% == BibLateX quality report for filippou2004fedeaslab:
% Missing required field 'type'
% ? Title looks like it was stored in title-case in Zotero

@book{fischer1999dualnumber,
  title = {Dual-Number Methods in Kinematics, Statics, and Dynamics},
  author = {Fischer, Ian S.},
  date = {1999},
  publisher = {CRC Press},
  location = {Boca Raton, Fla},
  isbn = {978-0-8493-9115-6},
  pagetotal = {219}
}
% == BibLateX quality report for fischer1999dualnumber:
% ? unused Call number ("TA352 .F57 1999")
% ? unused Library catalog ("Library of Congress ISBN")

@article{frostig2018compiling,
  title = {Compiling Machine Learning Programs via High-Level Tracing},
  author = {Frostig, Roy and Johnson, Matthew James and Leary, Chris},
  date = {2018},
  url = {https://cs.stanford.edu/~rfrostig/pubs/jax-mlsys2018.pdf},
  abstract = {We describe JAX, a domain-specific tracing JIT compiler for generating high-performance accelerator code from pure Python and Numpy machine learning programs. JAX uses the XLA compiler infrastructure to generate optimized code for the program subroutines that are most favorable for acceleration, and these optimized subroutines can be called and orchestrated by arbitrary Python. Because the system is fully compatible with Autograd, it allows forward- and reverse-mode automatic differentiation of Python functions to arbitrary order. Because JAX supports structured control flow, it can generate code for sophisticated machine learning algorithms while maintaining high performance. We show that by combining JAX with Autograd and Numpy we get an easily programmable and highly performant ML system that targets CPUs, GPUs, and TPUs, capable of scaling to multi-core Cloud TPUs.},
  langid = {english}
}
% == BibLateX quality report for frostig2018compiling:
% Missing required field 'journaltitle'
% ? unused Library catalog ("Zotero")

@inproceedings{geuzaine2012onelab,
  title = {{{ONELAB}}: {{Open Numerical Engineering LABoratory}}},
  shorttitle = {{{ONELAB}}},
  booktitle = {Proceedings of the 7th {{European Conference}} on {{Numerical Methods}} in {{Electromagnetism}} ({{NUMELEC2012}})},
  author = {Geuzaine, Christophe and Henrotte, François and Marchandise, Emilie and Remacle, Jean-François and Dular, Patrick and Sabariego, Ruth},
  date = {2012-01-01},
  abstract = {We present the ONELAB software, a lightweight open source toolkit to interface finite element and related solvers used in a variety of engineering disciplines, and to construct multi-code models with maximum flexibility, efficiency and user-friendliness.}
}
% == BibLateX quality report for geuzaine2012onelab:
% ? Title looks like it was stored in title-case in Zotero
% ? unused Library catalog ("ResearchGate")

@article{gilbert1992automatic,
  title = {Automatic Differentiation and Iterative Processes},
  author = {Gilbert, Jean Charles},
  date = {1992-01},
  journaltitle = {Optimization Methods and Software},
  shortjournal = {Optimization Methods and Software},
  volume = {1},
  number = {1},
  pages = {13--21},
  issn = {1055-6788, 1029-4937},
  doi = {10.1080/10556789208805503},
  abstract = {We identify a class of iterative processes that can be used in the definition of a function while preserving the good behavior of automatic differentiation codes on this function. By iterative process, we mean a process where the number of iterations is determined by a stopping criterion, which can depend on the independent variables. By good behavior, we mean that the derivatives will be calculated correctly. asymptotically. The class contams the Newton method.},
  langid = {english},
  annotation = {43 citations (Crossref) [2022-06-27]}
}
% == BibLateX quality report for gilbert1992automatic:
% 'issn': not a valid ISSN
% ? unused Library catalog ("DOI.org (Crossref)")
% ? unused Url ("https://www.tandfonline.com/doi/full/10.1080/10556789208805503")

@software{GMSH,
  title = {Gmsh},
  author = {Geuzaine, C. and Remacle, J. -F.},
  date = {2020},
  location = {Liege, Belgium.},
  url = {http://gmsh.info/},
  abstract = {GMSH  is a free 3D finite element mesh generator with a built-in CAD engine and post-processor. Its design goal is to provide a fast, light and user-friendly meshing tool with parametric input and advanced visualization capabilities. Gmsh is built around four modules: geometry, mesh, solver and post-processing. The specification of any input to these modules is done either interactively using the graphical user interface, in ASCII text files using Gmsh's own scripting language (.geo files), or using the C++, C, Python or Julia API.},
  version = {2.0}
}
% == BibLateX quality report for GMSH:
% ? unused Programming language ("Python, C++")

@online{google2017xla,
  title = {{{XLA}} - {{TensorFlow}}, Compiled},
  author = {{google}},
  date = {2017-03-06},
  url = {https://developers.googleblog.com/2017/03/xla-tensorflow-compiled.html},
  urldate = {2021-04-18},
  abstract = {News and insights on Google platforms, tools, and events.},
  langid = {english},
  organization = {Google Developers Blog}
}

@inproceedings{griewank1989automatic,
  title = {On {{Automatic Differentiation}}},
  booktitle = {In {{Mathematical Programming}}: {{Recent Developments}} and {{Applications}}},
  author = {Griewank, Andreas},
  date = {1989},
  pages = {83--108},
  publisher = {Kluwer Academic Publishers},
  abstract = {In comparison to symbolic differentiation and numerical differencing, the chain rule based technique of automatic differentiation is shown to evaluate partial derivatives accurately and cheaply. In particular it is demonstrated that the reverse mode of automatic differentiation yields any gradient vector at no more than five times the cost of evaluating the underlying scalar function. After developing the basic mathematics we describe several software implementations and briefly discuss the ramifications for optimization.}
}
% == BibLateX quality report for griewank1989automatic:
% ? Title looks like it was stored in title-case in Zotero
% ? unused Library catalog ("CiteSeer")

@article{griewank1993derivative,
  title = {Derivative Convergence for Iterative Equation Solvers},
  author = {Griewank, Andreas and Bischof, Christian and Corliss, George and Carle, Alan and Williamson, Karen},
  date = {1993-01},
  journaltitle = {Optimization Methods and Software},
  shortjournal = {Optimization Methods and Software},
  volume = {2},
  number = {3-4},
  pages = {321--355},
  issn = {1055-6788, 1029-4937},
  doi = {10.1080/10556789308805549},
  abstract = {When nonlinear equation solvers are applied to parameter-dependent problems, th r iterates can be interpreted as functions of these variable parameters. The derivatives (if they exist sf these iterated functions can be recursively evaluated by the forward mode of automatic differentiation 'hen one may ask whether and how fast these derivative \&dues converge to the derivative of the implici\textasciitilde{} olution function, which may be needed for parameter identification, sensitivity studies, or design op;imi\&ition. It is shown here that derivative convergence is achieved with an R-linear or possibly R-superlinear rate for a large class of memoryless contractions or secant updating methods. For a wide' class of multistep contractions, we obtain R-linear convergence of a simplified derivative recurrence, wh\$h is more economical and can be easily generalized to higher-order derivatives. We also formulate a c nstructive criterion for derivative convergence based on the implicit function theorem. All theoretical r e s j t i are confirmed by numerical experiments on small test examples.},
  langid = {english},
  annotation = {26 citations (Crossref) [2022-06-27]}
}
% == BibLateX quality report for griewank1993derivative:
% 'issn': not a valid ISSN
% ? unused Library catalog ("DOI.org (Crossref)")
% ? unused Url ("http://www.tandfonline.com/doi/abs/10.1080/10556789308805549")

@article{griewank2003mathematical,
  title = {A Mathematical View of Automatic Differentiation},
  author = {Griewank, Andreas},
  date = {2003-05},
  journaltitle = {Acta Numerica},
  volume = {12},
  pages = {321--398},
  publisher = {Cambridge University Press},
  issn = {1474-0508, 0962-4929},
  doi = {10/fvwgj5},
  abstract = {Automatic, or algorithmic, differentiation addresses the need for the accurate and efficient calculation of derivative values in scientific computing. To this end procedural programs for the evaluation of problem-specific functions are transformed into programs that also compute the required derivative values at the same numerical arguments in floating point arithmetic. Disregarding many important implementation issues, we examine in this article complexity bounds and other more mathematical aspects of the program transformation task sketched above.},
  langid = {english}
}
% == BibLateX quality report for griewank2003mathematical:
% Unexpected field 'publisher'
% 'issn': not a valid ISSN
% ? unused Url ("https://www.cambridge.org/core/journals/acta-numerica/article/mathematical-view-of-automatic-differentiation/B798FD0455A69D3B13FC28F60D8DF387")

@article{haukaas2006strategies,
  title = {Strategies for Finding the Design Point in Non-Linear Finite Element Reliability Analysis},
  author = {Haukaas, Terje and Der Kiureghian, Armen},
  date = {2006-04},
  journaltitle = {Probabilistic Engineering Mechanics},
  shortjournal = {Probabilistic Engineering Mechanics},
  volume = {21},
  number = {2},
  pages = {133--147},
  issn = {02668920},
  doi = {10.1016/j.probengmech.2005.07.005},
  langid = {english},
  annotation = {61 citations (Crossref) [2022-06-27]}
}
% == BibLateX quality report for haukaas2006strategies:
% ? unused Library catalog ("DOI.org (Crossref)")
% ? unused Url ("https://linkinghub.elsevier.com/retrieve/pii/S0266892005000469")

@article{hoffmann2016hitchhiker,
  title = {A {{Hitchhiker}}'s {{Guide}} to {{Automatic Differentiation}}},
  author = {Hoffmann, Philipp H. W.},
  date = {2016-07},
  journaltitle = {Numerical Algorithms},
  shortjournal = {Numer Algor},
  volume = {72},
  number = {3},
  eprint = {1411.0583},
  eprinttype = {arXiv},
  pages = {775--811},
  issn = {1017-1398, 1572-9265},
  doi = {10/ggsd72},
  abstract = {This article provides an overview of some of the mathematical principles of Automatic Differentiation (AD). In particular, we summarise different descriptions of the Forward Mode of AD, like the matrix-vector product based approach, the idea of lifting functions to the algebra of dual numbers, the method of Taylor series expansion on dual numbers and the application of the push-forward operator, and explain why they all reduce to the same actual chain of computations. We further give a short mathematical description of some methods of higher-order Forward AD and, at the end of this paper, briefly describe the Reverse Mode of Automatic Differentiation.},
  annotation = {14 citations (Crossref) [2022-06-27]}
}
% == BibLateX quality report for hoffmann2016hitchhiker:
% 'issn': not a valid ISSN
% ? Title looks like it was stored in title-case in Zotero
% ? unused Url ("http://arxiv.org/abs/1411.0583")

@book{hofmann2007conference,
  title = {Conference Record of {{POPL}} 2007: The 34th {{ACM SIGPLAN-SIGACT Symposium}} on {{Principles}} of {{Programming Languages}} : Papers Presented at the {{Symposium}}, {{January}} 17-19, 2007, {{French Riviera}}, {{Nice}}, {{France}}.},
  shorttitle = {Conference Record of {{POPL}} 2007},
  author = {Hofmann, Martin and {ACM Symposium on Principles of Programming Languages} and {ACM Symposium on Principles of Programming Languages} and {Association for Computing Machinery} and {ACM Special Interest Group on Programming Languages}},
  date = {2007},
  isbn = {978-1-59593-575-5},
  langid = {english},
  annotation = {OCLC: 1164890122}
}
% == BibLateX quality report for hofmann2007conference:
% ? unused Library catalog ("Open WorldCat")

@incollection{iri2009automatic,
  title = {Automatic Differentiation: Introduction, History and Rounding Error {{estimationAutomatic Differentiation}}: {{Introduction}}, {{History}} and {{Rounding Error Estimation}}},
  shorttitle = {Automatic Differentiation},
  booktitle = {Encyclopedia of {{Optimization}}},
  author = {Iri, Masao and Kubota, Koichi},
  editor = {Floudas, Christodoulos A. and Pardalos, Panos M.},
  date = {2009},
  pages = {153--159},
  publisher = {Springer US},
  location = {Boston, MA},
  doi = {10.1007/978-0-387-74759-0_26},
  abstract = {Article OutlineKeywordsIntroduction Algorithms ComplexityHistoryEstimates of Rounding ErrorsSee alsoReferences},
  isbn = {978-0-387-74759-0},
  langid = {english}
}
% == BibLateX quality report for iri2009automatic:
% ? unused Library catalog ("Springer Link")
% ? unused Url ("https://doi.org/10.1007/978-0-387-74759-0_26")

@software{JAX,
  title = {{{JAX}}},
  date = {2021-05-18T09:38:32Z},
  origdate = {2018-10-25T21:25:02Z},
  url = {https://github.com/google/jax},
  urldate = {2021-05-18},
  abstract = {Composable transformations of Python+NumPy programs: differentiate, vectorize, JIT to GPU/TPU, and more},
  organization = {Google}
}
% == BibLateX quality report for JAX:
% ? unused Library catalog ("GitHub")
% ? unused Programming language ("Python")

@software{Kratos,
  title = {Kratos {{Multiphysics}}},
  author = {Dadvand, Pooyan and Rossi, Riccardo and Oñate, Eugenio},
  date = {2010},
  origdate = {2017-02-13T10:58:24Z},
  location = {Barcelona, Spain},
  url = {http://www.cimne.com/kratos/},
  urldate = {2019-12-02},
  abstract = {Kratos Multiphysics (A.K.A Kratos) is a framework for building parallel multi-disciplinary simulation software. Modularity, extensibility and HPC are the main objectives. Kratos has BSD license and is written in C++ with extensive Python interface.},
  organization = {KratosMultiphysics},
  annotation = {github: https://github.com/KratosMultiphysics/Kratos}
}
% == BibLateX quality report for Kratos:
% ? Title looks like it was stored in title-case in Zotero
% ? unused Library catalog ("GitHub")
% ? unused Programming language ("C++")

@inproceedings{lattner2004llvm,
  title = {{{LLVM}}: A Compilation Framework for Lifelong Program Analysis Transformation},
  shorttitle = {{{LLVM}}},
  booktitle = {International {{Symposium}} on {{Code Generation}} and {{Optimization}}, 2004. {{CGO}} 2004.},
  author = {Lattner, C. and Adve, V.},
  date = {2004-03},
  pages = {75--86},
  doi = {10/d5brsd},
  abstract = {We describe LLVM (low level virtual machine), a compiler framework designed to support transparent, lifelong program analysis and transformation for arbitrary programs, by providing high-level information to compiler transformations at compile-time, link-time, run-time, and in idle time between runs. LLVM defines a common, low-level code representation in static single assignment (SSA) form, with several novel features: a simple, language-independent type-system that exposes the primitives commonly used to implement high-level language features; an instruction for typed address arithmetic; and a simple mechanism that can be used to implement the exception handling features of high-level languages (and setjmp/longjmp in C) uniformly and efficiently. The LLVM compiler framework and code representation together provide a combination of key capabilities that are important for practical, lifelong analysis and transformation of programs. To our knowledge, no existing compilation approach provides all these capabilities. We describe the design of the LLVM representation and compiler framework, and evaluate the design in three ways: (a) the size and effectiveness of the representation, including the type information it provides; (b) compiler performance for several interprocedural problems; and (c) illustrative examples of the benefits LLVM provides for several challenging compiler problems.},
  eventtitle = {International {{Symposium}} on {{Code Generation}} and {{Optimization}}, 2004. {{CGO}} 2004.}
}
% == BibLateX quality report for lattner2004llvm:
% ? Unsure about the formatting of the booktitle
% ? unused Library catalog ("IEEE Xplore")

@inproceedings{martins2001connection,
  title = {The Connection between the Complex-Step Derivative Approximation and Algorithmic Differentiation},
  booktitle = {39th {{Aerospace Sciences Meeting}} and {{Exhibit}}},
  author = {Martins, J. and Sturdza, Peter and Alonso, Juan},
  date = {2001-01-08},
  publisher = {{American Institute of Aeronautics and Astronautics}},
  location = {Reno,NV,U.S.A.},
  doi = {10.2514/6.2001-921},
  abstract = {This paper presents improvements to the complexstep derivative approximation method which increase its accuracy and robustness. These improvements unveil the connection to algorithmic differentiation theory. The choice between these two methods then hinges on a trade-off between ease of implementation and execution efficiency. Automatic implementations for Fortran and C/C++ are presented and their relative merits are discussed. These new methods were successfully implemented in two very large multidisciplinary programs and the resulting sensitivities are shown to be as accurate as the analyses. Accuracy and ease of implementation make these tools very attractive options for sensitivity analysis.},
  eventtitle = {39th {{Aerospace Sciences Meeting}} and {{Exhibit}}},
  langid = {english}
}
% == BibLateX quality report for martins2001connection:
% ? Unsure about the formatting of the booktitle
% ? unused Library catalog ("DOI.org (Crossref)")
% ? unused Url ("http://arc.aiaa.org/doi/10.2514/6.2001-921")

@article{mckenna2010nonlinear,
  title = {Nonlinear {{Finite-Element Analysis Software Architecture Using Object Composition}}},
  author = {McKenna, Frank and Scott, Michael H. and Fenves, Gregory L.},
  date = {2010-01},
  journaltitle = {Journal of Computing in Civil Engineering},
  shortjournal = {J. Comput. Civ. Eng.},
  volume = {24},
  number = {1},
  pages = {95--107},
  issn = {0887-3801, 1943-5487},
  doi = {10.1061/(ASCE)CP.1943-5487.0000002},
  abstract = {Object composition offers significant advantages over class inheritance to develop a flexible software architecture for finiteelement analysis. Using this approach, separate classes encapsulate fundamental finite-element algorithms and interoperate to form and solve the governing nonlinear equations. Communication between objects in the analysis composition is established using software design patterns. Root-finding algorithms, time integration methods, constraint handlers, linear equation solvers, and degree of freedom numberers are implemented as interchangeable components using the Strategy pattern. The Bridge and Factory Method patterns allow objects of the finite-element model to vary independently from objects that implement the numerical solution procedures. The Adapter and Iterator patterns permit equations to be assembled entirely through abstract interfaces that do not expose either the storage of objects in the analysis model or the computational details of the time integration method. Sequence diagrams document the interoperability of the analysis classes for solving nonlinear finite-element equations, demonstrating that object composition with design patterns provides a general approach to developing and refactoring nonlinear finite-element software.},
  langid = {english},
  annotation = {333 citations (Crossref) [2022-06-27]}
}
% == BibLateX quality report for mckenna2010nonlinear:
% 'issn': not a valid ISSN
% ? Title looks like it was stored in title-case in Zotero
% ? unused Library catalog ("DOI.org (Crossref)")
% ? unused Url ("http://ascelibrary.org/doi/10.1061/%28ASCE%29CP.1943-5487.0000002")

@article{mckenna2011opensees,
  title = {{{OpenSees}}: {{A Framework}} for {{Earthquake Engineering Simulation}}},
  shorttitle = {{{OpenSees}}},
  author = {McKenna, Frank},
  date = {2011-07},
  journaltitle = {Computing in Science \& Engineering},
  shortjournal = {Comput. Sci. Eng.},
  volume = {13},
  number = {4},
  pages = {58--66},
  issn = {1521-9615},
  doi = {10.1109/MCSE.2011.66},
  langid = {english},
  annotation = {408 citations (Crossref) [2022-06-27]}
}
% == BibLateX quality report for mckenna2011opensees:
% ? Title looks like it was stored in title-case in Zotero
% ? unused Library catalog ("DOI.org (Crossref)")
% ? unused Url ("http://ieeexplore.ieee.org/document/5931487/")

@software{MOOSE,
  ids = {2021idaholab},
  title = {Idaholab/Moose},
  date = {2021-02-03T14:41:55Z},
  origdate = {2014-02-14T22:26:15Z},
  location = {Idaho},
  url = {https://mooseframework.inl.gov/},
  urldate = {2021-04-18},
  abstract = {Multiphysics Object Oriented Simulation Environment},
  organization = {Idaho National Laboratory},
  annotation = {github: https://github.com/idaholab/moose}
}
% == BibLateX quality report for MOOSE:
% ? Title looks like it was stored in lower-case in Zotero
% ? unused Library catalog ("GitHub")
% ? unused Programming language ("C++")

@software{ONELAB,
  title = {{{ONELAB}}: {{Open Numerical Engineering LABoratory}}},
  shorttitle = {{{ONELAB}}},
  url = {http://onelab.info/},
  urldate = {2020-04-04},
  abstract = {ONELAB is an open-source, lightweight interface to finite element software. It is completely free: the default ONELAB software bundle contains the mesh generator Gmsh, the finite element solver GetDP and the optimization library conveks. Many other codes (free or not) can be easily interfaced as well.}
}
% == BibLateX quality report for ONELAB:
% ? Title looks like it was stored in title-case in Zotero

@software{OOFEM,
  title = {Oofem},
  date = {1999},
  origdate = {2017-10-31T16:02:29Z},
  location = {Prague, Czech Republic},
  url = {https://github.com/oofem/oofem},
  urldate = {2019-12-02},
  abstract = {OOFEM is free finite element code with object oriented architecture for solving mechanical, transport, fluid and multiphysics problems.},
  organization = {OOFEM}
}
% == BibLateX quality report for OOFEM:
% ? Title looks like it was stored in lower-case in Zotero
% ? unused Library catalog ("GitHub")
% ? unused Programming language ("C/C++")

@software{OpenSees,
  ids = {mckenna1999opensees},
  title = {{{OpenSees}}},
  author = {McKenna, Frank},
  date = {1999},
  origdate = {2018-09-12T00:49:34Z},
  location = {Berkeley, CA},
  url = {https://opensees.berkeley.edu/},
  urldate = {2019-06-26},
  abstract = {OpenSees Source Code Repository. Contribute to OpenSees/OpenSees development by creating an account on GitHub},
  key = {OpenSees},
  organization = {Pacific Earthquake Engineering Research Center}
}
% == BibLateX quality report for OpenSees:
% Unexpected field 'key'
% ? unused Library catalog ("GitHub")
% ? unused Programming language ("C++")

@article{ozaki1995higherorder,
  title = {Higher-Order Sensitivity Analysis of Finite Element Method by Automatic Differentiation},
  author = {Ozaki, I. and Kimura, F. and Berz, M.},
  date = {1995-07-01},
  journaltitle = {Computational Mechanics},
  shortjournal = {Computational Mechanics},
  volume = {16},
  number = {4},
  pages = {223--234},
  issn = {1432-0924},
  doi = {10/cqq274},
  abstract = {To design optimal mechanical structures, design sensitivity analysis technique using higher order derivatives are important. However, usual techniques for computing the derivatives, for example numerical differentiation methods, are hard to apply to real scale structures because of the large amount of computational time and the accumulation of computational errors.},
  langid = {english},
  annotation = {15 citations (Crossref) [2022-06-27]}
}
% == BibLateX quality report for ozaki1995higherorder:
% ? unused Library catalog ("Springer Link")
% ? unused Url ("https://doi.org/10.1007/BF00369867")

@inproceedings{pearlmutter2007lazy,
  title = {Lazy Multivariate Higher-Order Forward-Mode {{AD}}},
  booktitle = {Proceedings of the 34th Annual {{ACM SIGPLAN-SIGACT}} Symposium on {{Principles}} of Programming Languages  - {{POPL}} '07},
  author = {Pearlmutter, Barak A. and Siskind, Jeffrey Mark},
  date = {2007},
  pages = {155},
  publisher = {ACM Press},
  location = {Nice, France},
  doi = {10/bpjm8z},
  eventtitle = {The 34th Annual {{ACM SIGPLAN-SIGACT}} Symposium},
  isbn = {978-1-59593-575-5},
  langid = {english}
}
% == BibLateX quality report for pearlmutter2007lazy:
% ? Unsure about the formatting of the booktitle
% ? unused Library catalog ("DOI.org (Crossref)")
% ? unused Url ("http://portal.acm.org/citation.cfm?doid=1190216.1190242")

@article{permann2020moose,
  title = {{{MOOSE}}: {{Enabling}} Massively Parallel Multiphysics Simulation},
  shorttitle = {{{MOOSE}}},
  author = {Permann, Cody J. and Gaston, Derek R. and Andrš, David and Carlsen, Robert W. and Kong, Fande and Lindsay, Alexander D. and Miller, Jason M. and Peterson, John W. and Slaughter, Andrew E. and Stogner, Roy H. and Martineau, Richard C.},
  date = {2020-01-01},
  journaltitle = {SoftwareX},
  shortjournal = {SoftwareX},
  volume = {11},
  pages = {100430},
  issn = {2352-7110},
  doi = {10/ggsq4w},
  abstract = {Harnessing modern parallel computing resources to achieve complex multiphysics simulations is a daunting task. The Multiphysics Object Oriented Simulation Environment (MOOSE) aims to enable such development by providing simplified interfaces for specification of partial differential equations, boundary conditions, material properties, and all aspects of a simulation without the need to consider the parallel, adaptive, nonlinear, finite element solve that is handled internally. Through the use of interfaces and inheritance, each portion of a simulation becomes reusable and composable in a manner that allows disparate research groups to share code and create an ecosystem of growing capability that lowers the barrier for the creation of multiphysics simulation codes. Included within the framework is a unique capability for building multiscale, multiphysics simulations through simultaneous execution of multiple sub-applications with data transfers between the scales. Other capabilities include automatic differentiation, scaling to a large number of processors, hybrid parallelism, and mesh adaptivity. To date, MOOSE-based applications have been created in areas of science and engineering such as nuclear physics, geothermal science, magneto-hydrodynamics, seismic events, compressible and incompressible fluid flow, microstructure evolution, and advanced manufacturing processes.},
  langid = {english},
  annotation = {162 citations (Crossref) [2022-06-27]}
}
% == BibLateX quality report for permann2020moose:
% ? unused Library catalog ("ScienceDirect")
% ? unused Url ("https://www.sciencedirect.com/science/article/pii/S2352711019302973")

@software{pyhf,
  title = {Scikit-Hep/Pyhf},
  date = {2021-04-15T19:44:19Z},
  origdate = {2018-01-24T16:14:39Z},
  url = {https://github.com/scikit-hep/pyhf},
  urldate = {2021-04-17},
  abstract = {pure-Python HistFactory implementation with tensors and autodiff},
  organization = {Scikit-HEP Project}
}
% == BibLateX quality report for pyhf:
% ? Title looks like it was stored in lower-case in Zotero
% ? unused Library catalog ("GitHub")
% ? unused Programming language ("Python")

@software{sciann2021sciann,
  title = {Sciann/Sciann},
  author = {{sciann}},
  date = {2021-04-14T06:30:08Z},
  origdate = {2019-07-20T20:13:06Z},
  url = {https://github.com/sciann/sciann},
  urldate = {2021-04-17},
  abstract = {Deep learning for Engineers}
}
% == BibLateX quality report for sciann2021sciann:
% ? Title looks like it was stored in lower-case in Zotero
% ? unused Library catalog ("GitHub")
% ? unused Programming language ("Python")

@article{scott2008software,
  title = {Software {{Framework}} for {{Parameter Updating}} and {{Finite-Element Response Sensitivity Analysis}}},
  author = {Scott, Michael H. and Haukaas, Terje},
  date = {2008-09},
  journaltitle = {Journal of Computing in Civil Engineering},
  shortjournal = {J. Comput. Civ. Eng.},
  volume = {22},
  number = {5},
  pages = {281--291},
  issn = {0887-3801, 1943-5487},
  doi = {10.1061/(ASCE)0887-3801(2008)22:5(281)},
  abstract = {The finite-element software framework OpenSees is extended with parameter updating and response sensitivity capabilities to support client applications such as reliability, optimization, and system identification. Using software design patterns, member properties, applied loadings, and nodal coordinates can be identified and repeatedly updated in order to create customized finite-element model updating applications. Parameters are identified using a Chain of Responsibility software pattern, where objects in the finite-element model forward a parameterization request to component objects until the request is handled. All messages to identify and update parameters are passed through a Facade that decouples client applications from the finite-element domain of OpenSees. To support response sensitivity analysis, the Strategy design pattern facilitates multiple approaches to evaluate gradients of the structural response, whereas the Visitor pattern ensures that objects in the finite-element domain make the proper contributions to the equations that govern the response sensitivity. Examples demonstrate the software design and the steps taken by representative finite-element model updating and response sensitivity applications.},
  langid = {english},
  annotation = {19 citations (Crossref) [2022-06-27]}
}
% == BibLateX quality report for scott2008software:
% 'issn': not a valid ISSN
% ? Title looks like it was stored in title-case in Zotero
% ? unused Library catalog ("DOI.org (Crossref)")
% ? unused Url ("http://ascelibrary.org/doi/10.1061/%28ASCE%290887-3801%282008%2922%3A5%28281%29")

@article{spacone1996fibre,
  title = {Fibre {{Beam}}–{{Column Model}} for {{Non-Linear Analysis}} of {{R}}/{{C Frames}}: {{Part I}}. {{Formulation}}},
  shorttitle = {Fibre {{Beam}}–{{Column Model}} for {{Non-Linear Analysis}} of {{R}}/{{C Frames}}},
  author = {Spacone, E. and Filippou, F. C. and Taucer, F. F.},
  date = {1996},
  journaltitle = {Earthquake Engineering \& Structural Dynamics},
  shortjournal = {Earthquake Engng. Struct. Dyn.},
  volume = {25},
  number = {7},
  pages = {711--725},
  issn = {1096-9845},
  doi = {10.1002/(SICI)1096-9845(199607)25:7<711::AID-EQE576>3.0.CO;2-9},
  abstract = {This paper presents a fibre beam–column element for the non-linear static and dynamic analysis of reinforced concrete frames. It is assumed that plane sections remain plane and normal to the longitudinal axis. The effects of shear and bond-slip are, thus, presently neglected. The non-linear hysteretic behaviour of the element derives from the constitutive relations of concrete and reinforcing steel fibres into which each section is divided. The element formulation is flexibility-based and relies on force interpolation functions that strictly satisfy the equilibrium of bending moments and axial force along the element. Since the element does not make use of displacement interpolation functions, an iterative algorithm is needed for the determination of the resisting forces during the element state determination. The proposed algorithm is accurate and stable, even in the presence of strength loss, and is, thus, capable of tracing very well the highly non-linear behaviour of R/C members under cyclic load combinations of bending moment and axial force.},
  langid = {english},
  annotation = {570 citations (Crossref) [2022-06-27]\\
804 citations (Semantic Scholar/DOI) [2022-06-13]}
}
% == BibLateX quality report for spacone1996fibre:
% ? Title looks like it was stored in title-case in Zotero
% ? unused extra: _eprint ("https://onlinelibrary.wiley.com/doi/pdf/10.1002/%28SICI%291096-9845%28199607%2925%3A7%3C711%3A%3AAID-EQE576%3E3.0.CO%3B2-9")
% ? unused Library catalog ("Wiley Online Library")
% ? unused Url ("https://onlinelibrary.wiley.com/doi/abs/10.1002/%28SICI%291096-9845%28199607%2925%3A7%3C711%3A%3AAID-EQE576%3E3.0.CO%3B2-9")

@article{taylor2017feap,
  title = {{{FEAP}} - - {{Programmer Manual}}},
  author = {Taylor, Robert L and Govindjee, Sanjay},
  date = {2017},
  pages = {107},
  langid = {english}
}
% == BibLateX quality report for taylor2017feap:
% Missing required field 'journaltitle'
% ? Title looks like it was stored in title-case in Zotero
% ? unused Library catalog ("Zotero")

@article{tijskens2000automatic,
  title = {Automatic {{Differentiation}} for {{Solving Nonlinear Partial Differential Equations}}: {{An Efficient Operator Overloading Approach}}},
  author = {Tijskens, E and Roose, D and Ramon, H and Baerdemaeker, J De},
  date = {2000},
  pages = {43},
  abstract = {By resorting to Automatic Differentiation (AD) users of nonlinear PDE solvers can be relieved from the extra work of linearising a nonlinear PDE system and at the same time improve on the computational efficiency. This paper describes the main AD techniques and discusses how the operator overloading approach of AD can be extended to eliminate the overhead generally incurred with operator overloading. A recent AD system FastDer++, specially designed for this purpose, is integrated into a Least Squares solver. The necessary modifications to the general FEM algorithms. Code fragments and timing results demonstrate that (1) integrating AD with nonlinear PDE solvers leads to highly flexible code with a close resemblance to the mathematical expression of the problem, (2) coding and debugging efforts are greatly reduced, and (3) the computational efficiency is improved.},
  langid = {english}
}
% == BibLateX quality report for tijskens2000automatic:
% Missing required field 'journaltitle'
% ? Title looks like it was stored in title-case in Zotero
% ? unused Library catalog ("Zotero")

@article{vassilev2015clad,
  title = {Clad — {{Automatic Differentiation Using Clang}} and {{LLVM}}},
  author = {Vassilev, V and Vassilev, M and Penev, A and Moneta, L and Ilieva, V},
  date = {2015-05-22},
  journaltitle = {Journal of Physics: Conference Series},
  shortjournal = {J. Phys.: Conf. Ser.},
  volume = {608},
  pages = {012055},
  issn = {1742-6588, 1742-6596},
  doi = {10/gjqwtg},
  annotation = {4 citations (Crossref) [2022-06-27]}
}
% == BibLateX quality report for vassilev2015clad:
% 'issn': not a valid ISSN
% ? Title looks like it was stored in title-case in Zotero
% ? unused Library catalog ("DOI.org (Crossref)")
% ? unused Url ("https://iopscience.iop.org/article/10.1088/1742-6596/608/1/012055")

@article{walther2003adolc,
  title = {{{ADOL-C}}: {{Automatic Differentiation Using Operator Overloading}} in {{C}}++},
  shorttitle = {{{ADOL-C}}},
  author = {Walther, Andrea and Griewank, Andreas and Vogel, Olaf},
  date = {2003},
  journaltitle = {PAMM},
  volume = {2},
  number = {1},
  pages = {41--44},
  issn = {1617-7061},
  doi = {10/cbs2c2},
  abstract = {In this paper, we present two strategies for the implementation of Automatic Differentiation (AD) based on the operator overloading facility in C++. Subsequently, we describe the capabilities of the AD-tool ADOL-C that applies operator overloading to differentiate C- and C++-code. Finally, we discuss some applications of ADOL-C.},
  langid = {english},
  annotation = {20 citations (Crossref) [2022-06-27]}
}
% == BibLateX quality report for walther2003adolc:
% ? Title looks like it was stored in title-case in Zotero
% ? unused extra: _eprint ("https://onlinelibrary.wiley.com/doi/pdf/10.1002/pamm.200310011")
% ? unused Library catalog ("Wiley Online Library")
% ? unused Url ("https://onlinelibrary.wiley.com/doi/abs/10.1002/pamm.200310011")

@article{wengert1964simple,
  title = {A Simple Automatic Derivative Evaluation Program},
  author = {Wengert, R. E.},
  date = {1964-08-01},
  journaltitle = {Communications of the ACM},
  shortjournal = {Commun. ACM},
  volume = {7},
  number = {8},
  pages = {463--464},
  issn = {0001-0782},
  doi = {10/bwkd4g},
  abstract = {A procedure for automatic evaluation of total/partial derivatives of arbitrary algebraic functions is presented. The technique permits computation of numerical values of derivatives without developing analytical expressions for the derivatives. The key to the method is the decomposition of the given function, by introduction of intermediate variables, into a series of elementary functional steps. A library of elementary function subroutines is provided for the automatic evaluation and differentiation of these new variables. The final step in this process produces the desired function's derivative. The main feature of this approach is its simplicity. It can be used as a quick-reaction tool where the derivation of analytical derivatives is laborious and also as a debugging tool for programs which contain derivatives.},
  annotation = {156 citations (Crossref) [2022-06-27]}
}
% == BibLateX quality report for wengert1964simple:
% ? unused Library catalog ("Aug. 1964")
% ? unused Url ("https://doi.org/10.1145/355586.364791")

@report{wilson1970sap,
  title = {{{SAP}} - {{A General Structural Analysis Program}}},
  author = {Wilson, Edward L.},
  date = {1970-10},
  number = {SEMM-70-20},
  institution = {University of California Berkeley}
}
% == BibLateX quality report for wilson1970sap:
% Missing required field 'type'
% ? Title looks like it was stored in title-case in Zotero

@software{wilson1975etabs,
  title = {{{ETABS}}},
  author = {Wilson, Edward L},
  date = {1975},
  location = {Berkeley, Ca}
}

@unpublished{zwicke2017automatic,
  title = {Automatic Implementation of Material Laws: {{Jacobian}} Calculation in a Finite Element Code with {{TAPENADE}}},
  shorttitle = {Automatic Implementation of Material Laws},
  author = {Zwicke, Florian and Knechtges, Philipp and Behr, Marek and Elgeti, Stefanie},
  date = {2017-02-21},
  eprint = {1702.06365},
  eprinttype = {arXiv},
  eprintclass = {physics},
  doi = {10.1016/j.camwa.2016.10.010},
  langid = {english}
}